%
% mellin-integrand.tex -- 
%
% (c) 2021 Prof Dr Andreas Müller, OST Ostschweizer Fachhochschule
%
\documentclass[tikz]{standalone}
\usepackage{amsmath}
\usepackage{times}
\usepackage{txfonts}
\usepackage{pgfplots}
\usepackage{csvsimple}
\usetikzlibrary{arrows,intersections,math,calc}
\definecolor{darkred}{rgb}{0.8,0,0}
\definecolor{darkgreen}{rgb}{0,0.6,0}
\begin{document}
\def\skala{1.05}
\begin{tikzpicture}[>=latex,thick,scale=\skala]
		
	% Achsen
	\draw[->] (-0.2,0) -- (6.2,0)
	node[right] {$x$};
	
	\draw[->] (0,-0.2) -- (0,4.6);
	
	% Beschriftung der y-Achse bewusst links
	\node[anchor=east] at (-0.15,3.8)
	{$\displaystyle\frac{x^{\sigma-1}}{e^x-1}$};
	
	% Kurve: nahe genug an x=0 beginnen,
	% damit der singuläre Anstieg sichtbar wird
	\begin{scope}
		\clip (-0.2,-0.2) rectangle (6.2,4.6);
		\draw[thick,domain=0.08:6,samples=200]
		plot (\x,{(\x^0.4)/(exp(\x)-1)});
	\end{scope}
	
	% Hinweise
	\node[align=left,anchor=west] at (0.55,3.45)
	{Singularität\\bei $0$};
	
	\node[align=center] at (4.55,0.82)
	{exponentieller\\Abfall};
	
\end{tikzpicture}
\end{document}

